% BHU PYQ -- Manually transcribed & error-corrected
\documentclass[11pt,a4paper]{article}

\usepackage[a4paper,top=2.5cm,bottom=2.5cm,left=2.8cm,right=2.8cm,headheight=14pt]{geometry}
\usepackage{amsmath,amssymb}
\usepackage[T1]{fontenc}
\usepackage[utf8]{inputenc}
\usepackage{lmodern}
\usepackage{microtype}
\usepackage{fancyhdr}
\usepackage{enumitem}
\usepackage{hyperref}
\usepackage{lastpage}
\usepackage{parskip}

\hypersetup{colorlinks=true,urlcolor=black,linkcolor=black,
  pdftitle={BPT-603: Elements of Nuclear Physics -- BHU PYQ},pdfauthor={Banaras Hindu University}}

\pagestyle{fancy}\fancyhf{}
\renewcommand{\headrulewidth}{0.4pt}\renewcommand{\footrulewidth}{0.4pt}
\fancyhead[L]{\small   \textit{Banaras Hindu University}}
\fancyhead[R]{\small   \textit{BPT-603: Elements of Nuclear Physics}}
\fancyfoot[C]{\small Page  \thepage\ of \pageref{LastPage}}


\newlist{parts}{enumerate}{2}
\setlist[parts,1]{label=(\alph*),leftmargin=*,itemsep=5pt,topsep=3pt}
\setlist[parts,2]{label=(\roman*),leftmargin=*,itemsep=3pt,topsep=2pt}

\newcommand{\pts}[1]{\hfill   \textit{[#1]}}


\begin{document}


\begin{center}
  {\large\bfseries BANARAS HINDU UNIVERSITY}\\[2pt]
  {
\normalsize B.Sc. (Hons.) Semester VI Examination, 2013-14}\\[6pt]
  \rule{0.8\linewidth}{0.5pt}\\[6pt]
  {\large\bfseries Paper No. BPT-603: Elements of Nuclear Physics}\\[4pt]
  {
\normalsize Time: 3 Hours\hfill Maximum Marks: 70}
\end{center}
\rule{\linewidth}{0.4pt}
\smallskip



\noindent   \textit{Instructions: Answer all questions. Always draw necessary diagram or figure if needed. Symbols have their usual meanings.}
\bigskip




\noindent   \textbf{Question 1.} Answer any   \textbf{five} questions of the following: \pts{5 $ \times$ 3 = 15}

\begin{parts}
  \item Nucleon-nucleon forces are charge independent. Why do the $p-p$ and $n-n$ bound states not exist?
  \item On what factors does the choice of a good detector scintillator depend?
  \item Explain, with one example each, the following terms: (i) Mirror Nuclei, (ii) Singly and doubly magic nuclei.
  \item List one detector each which can be used for: (i) ionization and track measurements, (ii) time measurements, and (iii) particle identification.
  \item Explain why it is that only $\alpha$-particles are emitted by radioactive nuclei while protons and neutrons are not.
  \item A photon of energy $2.22\, \text{MeV}$ can disintegrate the deuteron into its components. Given that the masses of the neutron and proton are $1.008665\, \text{u}$ and $1.007277\, \text{u}$ respectively, what is the mass of the deuteron?
  \item What are the informations one can obtain from electron-nucleus scattering experiments?
\end{parts}

\medskip\rule{\linewidth}{0.2pt}




\noindent   \textbf{Question 2.}

\begin{parts}
  \item The radius of a nucleus is $R = r_0 A^{1/3}$ with $r_0 = 1.2\, \text{fm}$. Using Heisenberg's relations, estimate the mean kinetic energy and momentum of a nucleon inside a nucleus. \pts{5}
  \item Write down the Bethe-Weizs\"{a}cker mass formula and explain the various terms that contribute to it. \pts{4}
  \item Describe in brief the Fermi gas model of the nucleus. Show that the Fermi energy is independent of the number of nucleons in the nucleus for constant nuclear density. \pts{5}
\end{parts}

\medskip\rule{\linewidth}{0.2pt}




\noindent   \textbf{Question 3.}

\begin{parts}
  \item With the help of a neat diagram, describe the characteristics of the mode of operation of gas-filled detectors. \pts{5}
  \item List some of the important features of a scintillating signal and discuss at least two important features with suitable sketches. \pts{4}
  \item A $ \text{NaI}( \text{Tl})$ detector in the shape of a $5\, \text{cm}$ cube is bombarded by a beam of $2.8\, \text{MeV}$ gamma-radiation normal to one face of the cube. Calculate:
    
\begin{parts}
      \item What fraction of the gamma radiation is detected?
      \item What fraction of the detected gamma appears in the photopeak, Compton distribution, and pair peaks?
    \end{parts}
    Assume that there is no re-absorption of Compton gamma or annihilation quanta. (Given: attenuation coefficient $\mu$ for $2.8\, \text{MeV}$ gammas is $0.135\, \text{cm}^{-1}$, photoelectric coefficient $\mu_p = 2.5  \times 10^{-3}\, \text{cm}^{-1}$, Compton coefficient $\mu_{ \text{compton}} = 0.113\, \text{cm}^{-1}$, and pair production coefficient $\mu_{ \text{pair}} = 0.020\, \text{cm}^{-1}$). \pts{5}
\end{parts}

\medskip\rule{\linewidth}{0.2pt}




\noindent   \textbf{Question 4.}

\begin{parts}
  \item What are the informations we can obtain from the nucleon-nucleon ($n-p$ and $p-p$) scattering experiments? \pts{6}
  \item What do you know about the basic properties of the nucleus? Explain any two. Show that the density of nuclear matter is $\approx 10^{17}\, \text{kg/m}^3$ (or $\approx 10^5\, \text{tons/mm}^3$). (Assume that the mass of the proton is $1.67  \times 10^{-27}\, \text{kg}$ and that its spherical shape has a radius $R = r_0 A^{1/3}$, where $r_0 = 1.2  \times 10^{-13}\, \text{cm}$). \pts{7}
\end{parts}
\hfill   \textbf{OR}

\begin{parts}
  \item Define the $Q$-value of alpha decay. Write down the important features of alpha decay. What causes alpha decay? Explain it. \pts{6}
  \item Calculate the $Q$-value, kinetic energy $T_\alpha$, and the Coulomb barrier $V_c$ for the primary branch of the alpha decay of $^{212} \text{Po}$ to the ground state of $^{208} \text{Pb}$. Here, the mass excess of alpha is $2.4249\, \text{MeV}$, and excess masses of the daughter and parent are $-21.759\, \text{MeV}$ and $-10.381\, \text{MeV}$, respectively. \pts{7}
\end{parts}

\medskip\rule{\linewidth}{0.2pt}




\noindent   \textbf{Question 5.}

\begin{parts}
  \item What are the properties of the ground state of the deuteron? What are the informations about the nucleon-nucleon potential obtained from the experimental data on the deuteron? \pts{9}
  \item Make a neat sketch of the binding energy curve. Explain the basis for fission and fusion of nuclei. \pts{5}
\end{parts}
\hfill   \textbf{OR}

\begin{parts}
  \item What are elementary particles and how are they classified? Describe the different types of interactions that can occur between the elementary particles. \pts{7}
  \item Provide the names, masses, and charges of six quarks discovered so far. What do you understand by the colour degree of freedom? What is the need for introducing the colour degree of freedom for the quarks? \pts{7}
\end{parts}

\vfill
\rule{\linewidth}{0.4pt}

\begin{center}\small
  \href{https://bhu-exam.vercel.app/}{Click here to visit our website}\\[4pt]
  \href{https://me-aryan.vercel.app/}{Click here to visit our contributor}\\[4pt]
  \href{https://quashan-me.vercel.app/}{Click here to visit our contributor}
\end{center}

\end{document}
